\subsection{Родственные работы и место исследования}
\paragraph{Среды моделирования.}
iFogSim~\cite{ifogsim} предназначен для исследования управления ресурсами IoT/Fog, включая задержки, сеть и стоимость. iFogSim2~\cite{ifogsim2} расширяет его средствами мобильности, кластеризации и микросервисов. EdgeCloudSim~\cite{edgecloudsim} ориентирован на экспериментальную оценку краевых вычислений, а YAFS~\cite{yafs} даёт программируемое описание распределённых IoT-сценариев и размещения. Эти инструменты шире нашей модели. Новое ядро не объявляется их заменой: его назначение --- прозрачная проверка отдельных очередей, допущений и связи с внешним прогнозным интерфейсом. Сравнения быстродействия этих симуляторов здесь не проводились.

\paragraph{Размещение и прогноз.}
HEFT~\cite{heft} связывает ранги задач с оценками вычислений и коммуникаций; наш краткий пример использует заданное разбиение и внешний порядок устройств, а не воспроизводит HEFT. Обзор размещения приложений в тумане~\cite{fogsurvey} показывает разнообразие ограничений и целей. Для FEFC этот обзор задаёт контекст явно ограниченной постановки части 3: фиксированное разбиение связной задачи и отдельная проверка ресурсных, сетевых и временных ограничений полного плана. Непрерывное декларативное рассуждение FogBrain~\cite{fogbrain} относится к пересмотру размещения при изменении среды. Wang и соавторы~\cite{wang} используют прогноз будущих затрат; отличие нашего интерфейса состоит в отделении численного прогноза отклика от порядковых команд и последующей обязательной проверки полного плана. Гарантии алгоритмов с прогнозным советом, например работы Wei~\cite{advice}, зависят от конкретной задачи и не переносятся на наш критерий автоматически. Работа Lykouris и Vassilvitskii~\cite{lykouris} о кэшировании с обученными подсказками служит основанием разделять полезность прогноза и поведение алгоритма при его ошибках; доказанные для кэширования гарантии конкурентности не переносятся на размещение FEFC.

\paragraph{Предсказание задержки и неопределённость.}
Для V2X исследовались сочетания машинного обучения и статистического прогноза задержки на данных поездок~\cite{v2x}; для потоковых вычислений --- архитектуры с механизмом внимания~\cite{chu}. Их объекты, данные и цели отличаются от синтетического MLP данной работы, поэтому опубликованные выигрыши не приписываются нашей модели. Выбор метрики требует учитывать шкалу и цель прогноза~\cite{hyndman}; здесь MAE оценивается на одной и той же выборке в миллисекундах. Квантильный порог оценивается отдельно от точечной ошибки. Split-conformal методы~\cite{lei,angelopoulos} задают конечновыборочную поправку при обменности наблюдений. Очередной временной ряд зависим и может смещаться, поэтому показывается фактическое покрытие, а не неподтверждённая гарантия 95\%.

\paragraph{Очереди и ограничения сопоставления.}
Модели очередей различаются законами поступлений, обслуживания и дисциплиной~\cite{kendall}. Равномерный сервис рассматриваемого имитатора не совпадает с экспоненциальным сервисом M/M/1/K. Закон Литтла~\cite{little} используется с интенсивностью принятых, а не всех поступивших задач. Аналитическая кривая служит проверяемым приближением, а не эталоном точного совпадения. Вклад настоящей работы --- воспроизводимое описание и проверка связи этих компонентов; мировой приоритет, превосходство над перечисленными системами и промышленная готовность не заявляются.
